\documentclass[11pt,a4paper,fontset=none]{ctexart}
\usepackage[left=1.25cm,right=1.25cm,top=0.95cm,bottom=0.95cm]{geometry}
\usepackage{fontspec}
\usepackage{hyperref}
\usepackage{xcolor}
\usepackage{titlesec}
\usepackage{tabularx}
\usepackage{graphicx}
\usepackage{wrapfig}

\setmainfont{Times New Roman}
\setsansfont{Arial}
\setCJKmainfont{Noto Serif CJK SC}
\setCJKsansfont{Noto Sans CJK SC}

\definecolor{accent}{RGB}{15,78,110}
\definecolor{muted}{RGB}{78,78,78}
\hypersetup{colorlinks=true,urlcolor=accent,linkcolor=accent}
\pagestyle{empty}
\setlength{\parindent}{0pt}
\setlength{\parskip}{0pt}
\titleformat{\section}{\large\bfseries\color{accent}}{}{0pt}{}[\vspace{-2pt}\titlerule]
\titlespacing*{\section}{0pt}{5pt}{3.5pt}

\newcommand{\entry}[2]{%
  \begin{tabularx}{\textwidth}{@{}Xr@{}}\textbf{#1} & {\color{muted}#2}\end{tabularx}%
}
\newcommand{\meta}[1]{{\small\color{muted}#1}\par}
\newcommand{\projectdesc}[1]{\vspace{1pt}{\small #1\par}\vspace{3pt}}
\setlength{\columnsep}{10pt}
\setlength{\intextsep}{0pt}
\newcommand{\projectseparator}{%
  \WFclear
  \vspace{2pt}
  \noindent{\color{accent!70}\rule{\textwidth}{0.45pt}}\par
  \vspace{4pt}%
}

\begin{document}

\noindent
\begin{minipage}[t]{0.79\textwidth}
  \vspace{0pt}
  {\fontsize{18}{21}\selectfont\bfseries 任硕}\quad
  {\fontsize{13}{16}\selectfont Shuo REN}\par
  \vspace{3pt}
  {\small 邮箱：
    sren@cse.cuhk.edu.hk \,$\cdot$\,
    renshuo139@gmail.com \,$\cdot$\,
    主页：\href{https://shuoren-mike.github.io}{shuoren-mike.github.io}}\par
  \vspace{6pt}
  \textbf{香港中文大学计算机科学与工程系博士生，面向复杂工程流程研究长程、工具驱动的大语言模型智能体。} 研究工作涵盖多智能体物理设计优化、失败条件驱动的闭环修复与自动生成评测环境，以可验证的 EDA 和数学优化工具为真实任务基础。\par
\end{minipage}\hfill
\begin{minipage}[t]{0.17\textwidth}
  \vspace{0pt}
  \raggedleft
  \includegraphics[height=3.20cm,trim=22 18 22 6,clip]{../materials/assets/avatar/sren.pdf}
\end{minipage}\par

\vspace{3pt}

\section*{实习与教育经历}
\entry{Research Intern，香港科学园 EDA 研究实验室}{2026年5月--至今}
\projectdesc{开展 Agentic 3D-IC 研究，将 LLM 智能体与工业物理设计工具及反馈闭环结合，用于 3D-IC 划分与设计空间探索。构建覆盖任务分解、工具调用、结果校验与迭代优化的智能体工作流，以可验证的物理设计反馈支持长程规划、错误恢复与设计决策，并沉淀可复用的工具接口、执行轨迹与评测机制。}

\entry{ABCO 项目科研合作（CUHK 产学研合作项目）}{2025年3月--2026年4月}
\projectdesc{项目经费规模 100 万港币，获“优秀”评价（最高档）；构建 \href{https://arxiv.org/abs/2607.12788}{CLIP-3D} 背后的物理设计到体系结构反馈闭环，连接解析式 3D 布局规划与 gem5、McPAT、CACTI、HotSpot。}

\entry{香港中文大学（CUHK）}{2024年9月--至今}
\meta{计算机科学与工程博士，全额研究生奖学金；与\href{https://www.cse.cuhk.edu.hk/people/faculty/tsung-yi-ho/}{\textcolor{muted}{何宗易教授}}和\href{https://www.cse.cuhk.edu.hk/~byu/}{\textcolor{muted}{余备教授}}合作}

\entry{华中科技大学（HUST）}{2020--2024}
\meta{计算机科学与技术工学学士；GPA 3.93/4.00，专业前 1.5\%；\textbf{国家奖学金}}

\section*{代表性智能体研究}
\begin{wrapfigure}[6]{l}{0.285\textwidth}
  \vspace{0pt}
  \centering
  \includegraphics[width=0.27\textwidth]{../materials/assets/agent_projects/agenticpd_intro.pdf}
  \vspace{0pt}
\end{wrapfigure}
{\small\textbf{AgenticPD}（ASP-DAC 2027）\href{https://arxiv.org/abs/2607.04758}{[paper]} \href{https://github.com/shuoren-mike/AgenticPD}{[code]} 构建芯片物理设计多智能体优化系统；同等 20 次评估预算下，\textbf{六组设计/工艺组合的布线后 WNS 均优于 AutoTuner 与 ORFS-Agent}。设计跨阶段搜索、工具接口与状态复用机制，根据执行反馈选择优化阶段；在 DeepSeek、Qwen、Kimi 上的九组对比中，七组优于单智能体。\par}

\vspace{4pt}
\projectseparator

\begin{wrapfigure}[5]{l}{0.285\textwidth}
  \vspace{0pt}
  \centering
  \includegraphics[width=0.27\textwidth]{../materials/assets/agent_projects/agenticeco_intro.pdf}
  \vspace{0pt}
\end{wrapfigure}
{\small\textbf{AgenticECO} \href{https://arxiv.org/abs/2608.03738}{[paper]} \href{https://anonymous.4open.science/r/AgenticECO}{[code]} 构建芯片 3D-IC 自动修复系统，将缺陷诊断、并行候选探索与低扰动布线接入工具验证闭环。相同预算下，\textbf{九例缺陷中七例 DRC 清零，传统修复与全量重布线均为两例}；清零案例平均仅改动 0.66\% 的网络，不触碰时钟网络。跨模型复测九例全部清零，七例通过全部检查，平均网络改动率 0.44\%。\par}

\vspace{4pt}
\projectseparator

\begin{wrapfigure}[4]{l}{0.285\textwidth}
  \vspace{0pt}
  \centering
  \includegraphics[width=0.27\textwidth]{../materials/assets/agent_projects/autolpbench_intro.pdf}
  \vspace{0pt}
\end{wrapfigure}
{\small\textbf{A$^{2}$utoLPBench} \href{https://arxiv.org/abs/2607.02141}{[paper]} \href{https://anonymous.4open.science/r/AutoLPBench/}{[code]} 构建从自动出题、答案校验到智能体执行的评测平台，支持难度控制、Docker 一键评测及轨迹/token 成本分析。以逆 KKT 构造已知最优解，无需人工标注答案；评测六个模型系列，在 MAMO Complex 的 53 题子集上，\textbf{solver--critic 闭环使 Kimi-K2.5 解题率由 15.1\% 提升至 30.2\%}。\par}

\WFclear

\section*{EDA 研究背景}
{\small
\textbf{Partitioning-Free 3D-IC Floorplanning}（ASP-DAC 2026）\href{https://flians.github.io/pdf/3DFloorplan.pdf}{[paper]}：联合优化芯片模块布局与晶粒分配；五个 GSRC 基准上的线长均优于六种对比方法。\par
\vspace{1pt}
\textbf{CLIP-3D} \href{https://arxiv.org/abs/2607.12788}{[paper]}：打通芯片体系结构与物理设计评估；48 个热约束配置中，相对 shelf-pack 基线的实际 BIPS 平均提升 13.75\%，最高 24.77\%。作为 3D-IC 案例在 IDAS 2026 keynote 中介绍。\par
\vspace{1pt}
\textbf{Chiplet3D} \href{https://arxiv.org/abs/2607.09742}{[paper]}：联合优化芯粒引脚与 3D 热场，相对对比方法线长平均降低 39--43\%，峰值温度最多降低 45.9$^\circ$C；通过 3D-ICE 仿真验证温度。\par
}

\section*{技术能力}
{\small
\textbf{Agent/LLM：}重度使用 Claude Code 和 Codex；AI 原生开发、多智能体工具调用、结构化 API 与 Docker 评测。\par
\vspace{2pt}
\textbf{EDA/系统：}基于 OpenROAD/ORFS 开源 flow 的 end-to-end 物理设计能力；Innovus；STA 与 3D-IC 流程。\par
}

\end{document}
